\chapter{Introducción}
\label{ch:introduction}

% La introducción presenta el trabajo y prepara al lector; suele redactarse al
% final. Escribe un par de párrafos de presentación antes de la primera sección.

\section{Contexto}
% Encuadra el tema: ámbito, antecedentes y por qué es relevante actualmente.

\section{Problema}
% Describe con precisión el problema, la necesidad o la pregunta que aborda el TFG.

\section{Motivación}
% Razones (técnicas, personales, sociales) que justifican acometer el trabajo.

\section{Objetivos}
% Objetivo general y objetivos específicos, a ser posible medibles. Puedes
% enumerarlos con un entorno itemize/enumerate (ver chapters/00_examples.tex).

\section{Marco regulatorio}
% Normativa, estándares técnicos o aspectos legales aplicables al trabajo.

\section{Estructura del documento}
% Breve guía, capítulo a capítulo, de lo que el lector encontrará a continuación.
